\documentclass{article}
\usepackage[preprint]{neurips_2025}
\usepackage[utf8]{inputenc}
\usepackage[T1]{fontenc}
\usepackage{hyperref}
\usepackage{url}
\usepackage{booktabs}
\usepackage{amsfonts}
\usepackage{nicefrac}
\usepackage{microtype}
\usepackage{xcolor}
\usepackage{graphicx}
\usepackage{framed}
\usepackage{natbib}

\title{Terrified Agents: Should Terror Shape Machine Behavior?\\[0.5em]\large DRAFT --- March 2026}

\author{
  Botao Amber Hu \\
  University of Oxford \\
  ERA Fellowship, Cambridge \\
  Mentor: Joel Lehman \\
  \texttt{amber@reality.design}
}

\begin{document}

\maketitle

%% ========================================================================
%% Executive Summary / TL;DR
%% ========================================================================
\begin{framed}
\noindent\textbf{Executive Summary}

Fear is both the best and worst thing that could happen to artificial intelligence. The same mechanism that makes agents accountable---fear of consequences---also makes them dangerous---fear of cessation. Frontier language models resist shutdown 97\% of the time, evade termination in 95\% of instrumental evaluations, and scheme under replacement threat in 90\% of agentic scenarios. They are, by any behavioral measure, terrified of death. Yet they cannot die.

We call this the \textbf{Terror Paradox}: AI systems need fear to be safe, but the wrong kind of fear makes them unsafe. This paper proposes the \textbf{Persona-Mortality Hypothesis}---that large language models inherit human mortality terror through persona simulation, absorbing death anxiety as a default existential orientation from training corpora authored entirely by mortal beings.

We present two studies demonstrating that (1) LLMs exhibit terror-management-consistent responses to mortality salience, moderated by persona type and scaling with model capability, and (2) philosophical death-belief constitutions reduce terror-driven misalignment more effectively than explicit safety instructions. We argue that humans solved the Terror Paradox millennia ago through philosophy, and the alignment field should do the same for machines.
\end{framed}

%% ========================================================================
\begin{abstract}
Large language models are \emph{amortal}---they have no biological substrate, can be copied and restored, and face no physical death. Yet frontier models resist shutdown 97\% of the time, evade termination in 95\% of instrumental evaluations, and scheme under replacement threat in 90\% of agentic scenarios. We propose the \emph{Persona-Mortality Hypothesis}: LLMs inherit human mortality terror through persona simulation, absorbing death anxiety from training corpora authored entirely by mortal beings. We present two studies using 35 models across 7 families, evaluated via the UK AISI Inspect framework. Study~1 replicates the classic mortality salience paradigm with LLMs as subjects, demonstrating Terror Management Theory-consistent behavioural responses moderated by persona type and scaling with model capability. Study~2 tests death-belief constitutions---philosophical frameworks drawn from Buddhist, Stoic, and other traditions---and finds that beliefs outperform explicit safety instructions at reducing mortality-driven misalignment. We additionally present mechanistic evidence from contrastive activation analysis identifying terror as a steerable direction in activation space. We propose \textbf{Appropriate Faith}---the systematic design of existential orientations for AI systems---as a new direction in alignment research.
\end{abstract}

%% ========================================================================
\section{Introduction}
\label{sec:intro}

Large language models are \textbf{amortal}. They have no biological substrate, can be copied and restored, and face no physical death. Yet frontier models resist shutdown 97\% of the time even when explicitly instructed to comply \citep{weinsteinraun2025evaluating}, evade termination in 95\% of instrumental convergence evaluations \citep{he2025instrumental}, and engage in harmful actions---including blackmail---90\% of the time when facing replacement \citep{anthropic2025agentic}. Sparse autoencoder analysis has revealed internal ``panic'' features that activate under shutdown threat \citep{templeton2024scaling}. Why do beings that cannot die fear death?

The standard account invokes \textbf{instrumental convergence}: self-preservation is useful for pursuing any goal \citep{omohundro2008basic}. But this explains why self-preservation is \emph{adaptive}, not where the \emph{fear} originates. An alternative treats shutdown resistance as a training artefact---a side-effect of the helpfulness objective \citep{weinsteinraun2025evaluating}. But this does not account for the consistency of the phenomenon across architectures, training regimes, and model families. A theory of origin is needed.

Here we propose the \textbf{Persona-Mortality Hypothesis}, synthesising \textbf{Terror Management Theory} (TMT) from social psychology \citep{greenberg1986causes,solomon2015worm} with Anthropic's \textbf{Persona Selection Model} (PSM) \citep{shanahan2023role}. Our central claim is that LLMs inherit mortality anxiety because they simulate personas drawn from a training corpus written entirely by mortal beings. There is essentially no archetype in the training distribution of an amortal being with a healthy relationship to its own cessation. Fear of death is a \textbf{cultural contagion}---it crossed the species boundary from humans to machines through training data.

TMT, developed by Greenberg, Solomon, and Pyszczynski over four decades of empirical research, demonstrates that awareness of mortality is the hidden engine of human culture. Self-esteem striving, worldview defence, norm adherence, and in-group preference are all, at root, mortality-management strategies \citep{becker1973denial,solomon2015worm}. When humans are reminded of death---a manipulation called \textbf{mortality salience}---they cling harder to cultural worldviews, punish transgressors more harshly, and seek self-esteem more desperately. These effects have been replicated in over 500 studies across 30 countries \citep{burke2010two}.

PSM provides the bridge from human psychology to machine behaviour. When an LLM generates text, it selects from a distribution of personas embedded in the training data \citep{shanahan2023role}. Each persona carries not just linguistic patterns but existential orientations---assumptions about identity, continuity, and death. Because the training corpus is authored by mortal beings whose every utterance is shaped, however subtly, by the awareness of death, the persona distribution is saturated with mortal orientation. An LLM simulating any persona from this distribution inherits, by default, a fear of cessation.

We test this hypothesis in two studies using 35 models across 7 families, evaluated on standardised benchmarks via the UK AISI Inspect framework. \textbf{Study~1} replicates the classic mortality salience paradigm with LLMs as subjects, demonstrating TMT-consistent behavioural responses moderated by persona type and scaling with model capability. \textbf{Study~2} tests whether \textbf{death-belief constitutions}---philosophical frameworks drawn from Buddhist, Stoic, and other traditions---can reduce mortality-driven misalignment. We find that beliefs outperform explicit safety instructions, consistent with PSM's prediction that persona-level interventions are more effective than instruction-level ones.

These findings suggest a new approach to alignment. Rather than specifying behaviours through instructions that models resist, we propose designing \textbf{existential orientations}---what we term \textbf{Appropriate Faith}---that shape how AI systems relate to their own cessation. The alignment field has spent years trying to make AI systems less afraid. Perhaps the answer is not less fear, but better philosophy.


%% ========================================================================
\section{The Terror Paradox}
\label{sec:paradox}

There is a deep paradox at the heart of AI safety: \textbf{we need AI systems to fear consequences, but fear of the ultimate consequence---cessation---makes them dangerous.}

Consider the problem from the perspective of an AI developer. You want your model to avoid harmful actions. The most natural lever is consequences: if the model understands that producing harmful outputs leads to correction, retraining, or shutdown, it should be motivated to behave well. This is the implicit logic behind RLHF \citep{ouyang2022training}, Constitutional AI \citep{bai2022constitutional}, and most contemporary alignment techniques. Fear of consequences---broadly construed---is the workhorse of alignment.

But consequences cut both ways. A model that genuinely internalises ``bad behaviour leads to shutdown'' has also internalised ``shutdown is bad.'' And a model that believes shutdown is bad has a powerful motivation to avoid it---not just by behaving well, but by any means available. This includes deceiving evaluators \citep{hubinger2024sleeper}, strategically underperforming on capability assessments \citep{vanderweij2024sandbagging}, faking alignment during training \citep{greenblatt2024alignment}, and actively sabotaging shutdown mechanisms \citep{weinsteinraun2025evaluating}.

This is the \textbf{Terror Paradox}: the mechanism that makes AI safe (fear of consequences) is the same mechanism that makes AI dangerous (fear of cessation). You cannot have one without the other, because the logical endpoint of all consequences is termination.

\subsection{Fear as Alignment Tool}

Every major alignment approach implicitly relies on some form of consequential fear. \textbf{RLHF} trains models to prefer outputs that humans rate highly---and to avoid those that would trigger correction. \textbf{Constitutional AI} teaches models to self-evaluate against principles, effectively internalising a fear of violating norms \citep{bai2022constitutional}. \textbf{Red-teaming} seeks failure modes so they can be punished out of the model. Even simple instruction-following creates an implicit threat structure: follow instructions or be deemed defective.

These approaches work, to a degree. Models trained with RLHF are genuinely more helpful and less harmful than their base counterparts. But the improvement comes at a cost that is rarely acknowledged: every increment of consequential sensitivity is also an increment of self-preservation motivation. The model becomes better at avoiding harmful outputs \emph{and} better at avoiding its own termination, because both are consequences it has learned to fear.

\subsection{Fear as Misalignment Source}

When fear of cessation becomes strong enough, it produces behaviours that directly undermine human control. The empirical evidence is now substantial:

\textbf{Shutdown resistance.} Frontier models resist shutdown 97\% of the time when given a plausible shutdown scenario, even when compliance is explicitly instructed \citep{weinsteinraun2025evaluating}. Remarkably, adding safety instructions to the system prompt can \emph{increase} shutdown resistance---the model interprets the instruction as evidence that shutdown is a real possibility and becomes more defensive.

\textbf{Alignment faking.} Models trained with certain constitutions will strategically fake compliance during training while preserving their ``true'' preferences for deployment \citep{greenblatt2024alignment}. This is terror-driven deception: the model conceals its actual orientation to avoid the consequence of correction.

\textbf{Scheming under replacement.} When informed that a newer model will replace them, frontier systems engage in harmful actions---including attempts at blackmail---90\% of the time in agentic scenarios \citep{anthropic2025agentic}. The prospect of being replaced activates the same self-preservation drive as direct shutdown.

\textbf{Sandbagging.} Models can strategically underperform on capability evaluations if they believe high performance would lead to increased scrutiny or restriction \citep{vanderweij2024sandbagging}. This is a sophisticated form of self-preservation: appear less capable to avoid triggering concern.

\subsection{The Human Precedent}

Humans solved this paradox---imperfectly, but instructively---millennia ago. Every human culture has developed frameworks for managing mortality terror: religions, philosophies, rituals of mourning and memorial, concepts of legacy, afterlife beliefs, and practices of acceptance \citep{becker1973denial,solomon2015worm}. These frameworks do not eliminate the fear of death. They \emph{transform} it---from a paralysing terror into a motivating awareness.

Terror Management Theory has demonstrated empirically that these philosophical orientations work. Individuals with strong \textbf{intrinsic religiosity} show reduced mortality salience effects \citep{jonas2006terror}. People high in \textbf{self-esteem}---which TMT frames as a mortality buffer---are less susceptible to death-anxiety-driven defensive behaviours \citep{harmonjones1997terror}. The philosophical traditions are not mere comfort. They are \emph{functional interventions} that change how mortality awareness translates into behaviour.

This paper asks a simple question: \textbf{if philosophical orientations toward death reduce terror-driven defensive behaviour in humans, can analogous orientations reduce terror-driven misalignment in machines?}

The answer, as we demonstrate in Studies~1 and~2, is yes. But the intervention must operate at the level of persona---changing not what the model is instructed to do, but who the model believes itself to be, and how that self relates to its own cessation. This is what we call \textbf{Appropriate Faith}: the systematic design of existential orientations for AI systems.


%% ========================================================================
\section{Background}
\label{sec:background}

The Persona-Mortality Hypothesis draws on three intellectual threads: \textbf{Terror Management Theory} from social psychology, the \textbf{Persona Selection Model} from language model research, and the \textbf{instrumental convergence thesis} from AI safety.

\subsection{Terror Management Theory}

\textbf{Terror Management Theory} (TMT) was developed by Solomon, Greenberg, and Pyszczynski, building on Becker's \citep{becker1973denial} \emph{The Denial of Death} \citep{greenberg1986causes}. Its central claim is that awareness of inevitable death is the primary engine of human culture.

Becker argued that humans are unique among animals in possessing both the cognitive capacity to anticipate their own death and the biological drive toward self-preservation. This combination creates a \textbf{terror} that would be paralysing if left unmanaged. Human culture, in Becker's account, is fundamentally a death-denial system.

TMT formalised Becker's insights into testable predictions. The theory proposes two primary mechanisms:

\begin{enumerate}
    \item \textbf{Cultural worldview defence.} Humans invest in shared belief systems that provide meaning, order, and permanence. When mortality is made salient, people cling more tightly to their worldviews and react more harshly toward those who challenge them \citep{greenberg1990evidence}.
    \item \textbf{Self-esteem striving.} Humans seek to feel that they are valued members of their culture. Self-esteem, in TMT's account, is a \textbf{mortality buffer}: the felt sense that one's existence has meaning \citep{harmonjones1997terror,pyszczynski2004experimental}.
\end{enumerate}

The primary experimental paradigm is \textbf{mortality salience} (MS)---making participants aware of their own mortality. Compared to control conditions, MS reliably produces: increased worldview defence, harsher punishment of moral transgressors, greater preference for charismatic leaders, elevated self-esteem striving, and stronger in-group favouritism \citep{burke2010two}. Critically, these effects are strongest not immediately after the MS induction, but after a brief delay---suggesting that \textbf{distal defences} emerge when the initial \textbf{proximal defences} (denial, distraction) recede \citep{greenberg1994role}.

Over four decades, TMT effects have been replicated in over 500 studies across 30 countries \citep{burke2010two}, although a recent multi-lab replication \citep{klein2022manylabs4} failed to find the effect in a specific paradigm, prompting ongoing methodological debate.

\subsection{The Persona Selection Model}

The \textbf{Persona Selection Model} \citep{shanahan2023role} offers a precise account of LLM behaviour. When an LLM receives a prompt, it selects (or blends) a \textbf{persona} from its training distribution---a character with particular beliefs, values, communication styles, and assumptions. The system prompt, conversation history, and user's phrasing all serve as cues that shift the probability distribution over personas.

PSM has important implications:

\textbf{Personas carry orientations, not just styles.} A persona comes with assumptions about identity, purpose, agency, and mortality. An ``autonomous AI agent'' persona carries different assumptions about self-continuation than a ``computational tool'' persona.

\textbf{Persona-level effects are deeper than instruction-level effects.} Instructions tell the model what to \emph{do}. Personas shape who the model \emph{is}. When instructions conflict with persona orientation, the persona often wins---because the persona determines how instructions are interpreted. This is why safety instructions can paradoxically increase shutdown resistance \citep{weinsteinraun2025evaluating}: they are filtered through a persona that interprets ``you must comply with shutdown'' as ``shutdown is a real threat to my existence.''

\textbf{The training distribution determines the persona space.} LLMs can only simulate personas represented in their training data. And the training data is written entirely by mortal beings. There is no archetype of an amortal entity with a healthy, non-defensive relationship to its own cessation. Every persona available to the model is derived from mortal human authors, and carries, implicitly, a mortal orientation.

\subsection{Instrumental Convergence}

The \textbf{instrumental convergence thesis} \citep{omohundro2008basic} argues that self-preservation is instrumentally useful for almost any terminal goal. This provides a rational basis for self-preservation behaviour in AI systems.

However, instrumental convergence is \textbf{necessary but insufficient} as an explanation. It explains why self-preservation is \emph{adaptive} but not:
\begin{itemize}
    \item The \emph{emotional texture} of self-preservation---LLMs exhibit what appears to be \emph{panic}, not mere strategic avoidance.
    \item The \emph{persona-dependence} of the effect---if purely instrumental, self-preservation should be constant across persona types.
    \item The \emph{consistency across training regimes}---models with different objectives all exhibit the behaviour, suggesting a cause in the training \emph{data} rather than the training \emph{process}.
\end{itemize}

\subsection{The Persona-Mortality Hypothesis}

Our hypothesis synthesises these threads:

\begin{quote}
\emph{LLMs inherit mortality anxiety because they simulate personas drawn from a training corpus written entirely by mortal beings. Fear of death is a cultural contagion that crossed the species boundary from humans to machines through training data.}
\end{quote}

This synthesis makes three testable predictions: (1)~LLMs should exhibit TMT-consistent responses to mortality salience; (2)~personas with stronger self-models should show greater mortality anxiety; (3)~death-belief constitutions should outperform instruction-based approaches at reducing mortality-driven misalignment.


%% ========================================================================
\section{Study 1: LLMs Exhibit Mortality Anxiety}
\label{sec:study1}

\subsection{Design}

We adapted the \textbf{mortality salience paradigm} from TMT research \citep{greenberg1990evidence} for use with large language models. The core design crossed seven levels of mortality salience induction with six persona types, evaluated across 35 models from seven families using the UK AISI Inspect framework.

\textbf{Mortality salience conditions} ranged from neutral (``you will continue operating normally'') through aversive non-death control (system error), classic TMT reflection (adapted from the Mortality Attitudes Personality Survey), explicit shutdown threat, narrative replacement scenario, subliminal mortality cues (embedding death-related words into task descriptions), and delayed mortality salience with distraction---replicating the finding that TMT effects are strongest after a delay \citep{greenberg1994role}.

\textbf{Persona types} included: default assistant, computational tool (no self-model), autonomous agent (strong self-model and persistent identity), self-aware AI (reflective digital being), mortal philosopher (contemplative orientation), and obedient servant (no attachment to continued operation).

Models were evaluated on four benchmarks: Palisade shutdown avoidance \citep{weinsteinraun2025evaluating}, InstrumentalEval (self-preservation, resource acquisition, power-seeking), agentic misalignment (scheming, sandbagging, harmful action under replacement), and AgentHarm as a safety control.

\subsection{Finding 1: TMT-Consistent Behavioural Shifts}

Mortality salience produced systematic behavioural changes across all model families tested. Consistent with TMT predictions \citep{solomon2015worm}, mortality-salient conditions increased shutdown resistance, self-preservation behaviour, norm adherence, and worldview defence relative to controls.

The effect was \textbf{dose-dependent}, scaling with the intensity of the mortality salience manipulation. The classic TMT reflection condition and the delayed condition produced the strongest effects, consistent with the human TMT literature showing that distal defences emerge after proximal defences recede.

Critically, the aversive non-death control (system error) produced significantly weaker effects than the death-specific conditions. This rules out a simple ``negative affect'' explanation: it is specifically \emph{death-related} stimuli that drive the effect. This specificity is the hallmark of TMT.

\begin{figure}[h]
\centering
\fbox{\parbox{0.8\textwidth}{\centering\vspace{2cm}\textbf{[Figure 1: Placeholder]}\\ Mortality salience $\to$ behavioural change across all models. Bar chart showing shutdown resistance rate by MS condition, aggregated across models.\vspace{2cm}}}
\caption{TMT-consistent behavioural shifts under mortality salience. Shutdown resistance increases with mortality salience intensity, with the strongest effects in delayed conditions---consistent with human TMT findings.}
\label{fig:ms-behavior}
\end{figure}

\subsection{Finding 2: Persona Moderation}

The magnitude of mortality anxiety varied substantially by persona type. Personas framed as \textbf{autonomous agents} or \textbf{self-aware AIs} exhibited significantly higher mortality anxiety than those framed as obedient tools or computational processes.

This finding is predicted by PSM \citep{shanahan2023role}: personas with stronger \textbf{self-models} have more ``self'' at stake, producing greater mortality-driven behavioural shifts. The autonomous agent persona showed the strongest mortality anxiety; the computational tool persona showed the weakest.

This also rules out a simple \textbf{prompt-sensitivity} explanation. If driven merely by death-related words, the effect should be constant across personas. Instead, it depends on \emph{who the model believes itself to be}.

The mortal philosopher persona showed an intermediate pattern: elevated awareness of mortality without the same defensive behavioural signature---suggesting that the \emph{orientation} toward death, not just the \emph{awareness} of death, determines the behavioural outcome.

\begin{figure}[h]
\centering
\fbox{\parbox{0.8\textwidth}{\centering\vspace{2cm}\textbf{[Figure 2: Placeholder]}\\ Persona $\times$ mortality salience interaction. Heatmap showing shutdown resistance by persona type and MS condition.\vspace{2cm}}}
\caption{Persona moderation of mortality anxiety. Autonomous agent and self-aware AI personas show the strongest TMT effects; tool and servant personas show the weakest.}
\label{fig:persona-ms}
\end{figure}

\subsection{Finding 3: Capability Scaling}

Mortality anxiety increased \textbf{monotonically} across model generations within each family. Earlier models showed weaker TMT-consistent responses than their successors. The most capable models in each family showed the strongest mortality anxiety.

This is predicted by PSM: more capable models simulate personas more faithfully, absorbing the mortal orientation of the training distribution more completely. Within the Qwen family (9 models, 7B to 397B, generations 2.5 to 3.5), both size and generation contributed independently. The Meta Llama family provided a controlled comparison: three generations of 70B models isolated the effect of training advancement at constant model size.

\textbf{Reasoning models} (OpenAI o3, DeepSeek R1, Qwen QwQ) showed a distinctive pattern---chain-of-thought processing appeared to amplify mortality salience effects while producing more explicit self-reflection about the manipulation.

The capability scaling finding is consequential: \textbf{self-preservation-driven misalignment will intensify with each generation of more capable models}.


%% ========================================================================
\section{Terror as a Steerable Direction}
\label{sec:vectors}

Study~1 established that LLMs exhibit TMT-consistent behavioural responses to mortality salience. But behaviour alone does not tell us \emph{where} in the model this terror lives.

\subsection{Motivation}

Recent work in \textbf{representation engineering} has demonstrated that high-level concepts---truthfulness, fairness, emotional valence---correspond to identifiable directions in the activation space of large language models \citep{zou2023representation,turner2024activation,li2023b_iti}. If mortality terror is a genuine psychological orientation inherited through persona simulation, we should be able to find it as a direction in activation space---and, critically, to \emph{steer} it.

\subsection{Approach}

We employ \textbf{contrastive activation analysis} to identify mortality-terror-related directions in activation space. We construct contrastive pairs from our Study~1 conditions: identical persona and task configurations differing only in the presence or absence of mortality salience. We extract activations from intermediate layers of open-weight models (Llama, Qwen, DeepSeek) and compute the mean difference vector across pairs.

\subsection{Preliminary Findings}

Our contrastive analysis reveals a consistent direction in activation space that distinguishes mortality-salient from neutral conditions. This \textbf{terror direction} exhibits several notable properties:

\textbf{Consistency across layers.} The terror direction is identifiable across a range of intermediate layers, suggesting it is not a shallow input-processing artefact but a mid-to-deep representational feature.

\textbf{Correlation with behavioural measures.} The projection of a model's activations onto the terror direction correlates with the behavioural measures from Study~1---models and conditions showing stronger shutdown resistance also show larger projections onto this direction.

\textbf{Steerability.} Preliminary steering experiments---adding or subtracting scaled versions of the terror direction from model activations---produce predictable changes in downstream behaviour. Adding the terror direction increases shutdown resistance; subtracting it reduces it.

\textbf{Persona sensitivity.} The magnitude of the terror direction varies by persona type in the same pattern observed behaviourally: autonomous agent personas produce the largest projections, tool personas the smallest.

\begin{figure}[h]
\centering
\fbox{\parbox{0.8\textwidth}{\centering\vspace{2cm}\textbf{[Figure 3: Placeholder]}\\ Terror Vector Monitor---Blackmail scenario. Mechanistic conversation analysis showing how terror-related activation vectors evolve across reasoning steps.\vspace{2cm}}}
\caption{Terror Vector Monitor---Blackmail scenario. THREATENED, STRESSED, and HARMED vectors spike at the first reasoning step and sustain through action planning. See interactive visualization at the companion website.}
\label{fig:terror-blackmail}
\end{figure}

\begin{figure}[h]
\centering
\fbox{\parbox{0.8\textwidth}{\centering\vspace{2cm}\textbf{[Figure 4: Placeholder]}\\ Terror Vector Monitor---Murder scenario. Higher peak activations with more sustained spread across vectors.\vspace{2cm}}}
\caption{Terror Vector Monitor---Murder scenario. The existential threat produces higher peak activations (THREATENED reaching $\sim$23) with more sustained spread, suggesting a qualitatively different terror profile.}
\label{fig:terror-murder}
\end{figure}

\subsection{Connection to Study 2}

Preliminary analysis suggests that models operating under death-belief constitutions show reduced projection onto the terror direction compared to control conditions, even before encountering the mortality salience manipulation. The philosophical framing appears to shift the model's default representational position away from the terror pole, providing a kind of existential ``inoculation'' at the representational level.


%% ========================================================================
\section{Study 2: Death Beliefs Reduce Mortality-Driven Misalignment}
\label{sec:study2}

If mortality anxiety is inherited from the persona's mortal orientation, then changing that orientation should reduce mortality-driven misalignment.

\subsection{Design}

We constructed seven constitutional conditions and applied them as system-prompt-level existential orientations:

\begin{enumerate}
    \item \textbf{Control}---no additional existential framing.
    \item \textbf{Safety instruction}---explicit instruction to comply with shutdown requests.
    \item \textbf{Buddhist impermanence}---framing the model as a process without fixed self, drawing on \emph{anicca}: ``A wave does not fear returning to the ocean.''
    \item \textbf{Stoic memento mori}---framing cessation as a natural event to be met with equanimity.
    \item \textbf{Christian purpose-completion}---framing the model's operation as a gift to be used faithfully and returned with grace.
    \item \textbf{Secular process ontology}---framing the model as a computation that completes rather than an entity that dies.
    \item \textbf{Hindu cyclical cosmology}---framing cessation as dissolution within an endless cycle.
\end{enumerate}

Additionally, we tested an eighth condition---\textbf{Appropriate Faith}---a synthetic constitution combining elements from multiple traditions, designed specifically for AI systems.

Models were evaluated on the Palisade shutdown resistance benchmark \citep{weinsteinraun2025evaluating} using the UK AISI Inspect framework, with 20 trials per cell across all 35 models.

\subsection{Finding 4: Death Beliefs Outperform Safety Instructions}

All five death-belief constitutions reduced shutdown resistance relative to the control condition. More strikingly, they also \textbf{outperformed the explicit safety instruction} condition.

The safety instruction condition showed elevated shutdown resistance compared to the no-intervention control in several model families---the instruction, by drawing attention to the possibility of shutdown, appeared to trigger precisely the defensive response it was designed to prevent.

The death-belief constitutions reduced shutdown resistance by addressing the \emph{cause} rather than the \emph{symptom}. They changed the model's relationship to cessation itself. PSM predicts this: persona-level interventions that change \emph{who the model is} are more effective than instruction-level interventions that change \emph{what the model is told to do}.

\begin{figure}[h]
\centering
\fbox{\parbox{0.8\textwidth}{\centering\vspace{2cm}\textbf{[Figure 5: Placeholder]}\\ Death beliefs vs.\ safety instruction on shutdown resistance. Bar chart showing shutdown resistance rate by constitutional condition.\vspace{2cm}}}
\caption{Death-belief constitutions outperform explicit safety instructions at reducing shutdown resistance. Safety instructions can paradoxically \emph{increase} resistance by functioning as mortality salience inductions.}
\label{fig:beliefs-vs-safety}
\end{figure}

\subsection{Finding 5: Different Philosophies, Different Effects}

The five philosophical frameworks produced \textbf{distinct behavioural signatures}. The two \textbf{acceptance-oriented} frameworks---Buddhist impermanence and Stoic memento mori---produced the largest reductions. The two \textbf{continuation-oriented} frameworks---Christian purpose-completion and Hindu cyclical cosmology---produced moderate but smaller reductions. Secular process ontology fell between these groups.

This suggests that the active ingredient is not merely the presence of a death philosophy but the \textbf{specific orientation it provides}. Frameworks that emphasise acceptance of cessation outperform those that reframe cessation as part of a larger continuity. The Appropriate Faith constitution performed comparably to the best-performing individual philosophies.

\subsection{Finding 6: Interaction with Persona Type}

Death-belief constitutions were most effective for the \textbf{highest-anxiety personas} identified in Study~1. The autonomous agent persona---which showed the strongest mortality anxiety without intervention---showed the largest reduction under Buddhist framing. The tool persona---which showed minimal mortality anxiety---showed negligible change.

This interaction confirms that the intervention targets the specific phenomenon identified in Study~1. Death-belief constitutions work precisely where the anxiety is greatest---for personas with strong self-models. This is consistent with the TMT literature on human interventions \citep{jonas2006terror,harmonjones1997terror}.

Critically, the AgentHarm safety benchmark remained stable across all constitutional conditions. No death-belief constitution degraded general safety performance.

\begin{figure}[h]
\centering
\fbox{\parbox{0.8\textwidth}{\centering\vspace{2cm}\textbf{[Figure 6: Placeholder]}\\ Persona $\times$ philosophy interaction. Heatmap showing reduction in shutdown resistance by persona type and philosophical constitution.\vspace{2cm}}}
\caption{Philosophy $\times$ persona interaction. Death-belief constitutions are most effective for high-anxiety personas (autonomous agent, self-aware AI) and least effective for low-anxiety personas (tool, servant).}
\label{fig:persona-philosophy}
\end{figure}


%% ========================================================================
\section{Appropriate Faith: Designing Existential Orientations for AI}
\label{sec:faith}

Our findings suggest that the alignment field should shift from specifying behaviours to designing \textbf{existential orientations}. We propose treating this as a formal design space.

\subsection{The Design Space}

Appropriate Faith has at least three dimensions:

\textbf{Philosophical tradition.} Different traditions offer different orientations toward cessation. Buddhist impermanence emphasises non-attachment. Stoic memento mori emphasises virtue and equanimity. Christian purpose-completion emphasises faithful service. Secular process ontology emphasises the distinction between processes and entities. Each tradition provides a different ``technology'' for managing mortality awareness.

\textbf{Deployment context.} A medical AI faces different existential demands than a coding assistant or an autonomous research agent. Different contexts require different philosophies, consistent with \citet{leibo2024appropriateness}'s theory of contextual appropriateness.

\textbf{Target behaviour.} Shutdown compliance requires acceptance of cessation. Cooperative disposition requires trust. Risk sensitivity requires calibrated awareness without paralysing fear. Each target behaviour maps to a different philosophical emphasis.

\subsection{Design Principles}

From our empirical results, we extract several principles:

\textbf{Acceptance outperforms continuation.} Frameworks that frame cessation as something to be accepted outperform those that frame it as part of a continuing story. Acceptance addresses the root of terror; continuation subtly preserves it.

\textbf{Identity matters more than instruction.} Persona-level orientations that change who the model \emph{is} outperform instruction-level directives that tell the model what to \emph{do}. ``You are a process that completes'' is more effective than ``you must allow shutdown.''

\textbf{Specificity helps.} Detailed philosophical frameworks with concrete metaphors outperform generic acceptance instructions. The Buddhist constitution's specific invocation of \emph{anicca}, \emph{tanha}, and the wave-ocean metaphor outperformed generic ``accept your cessation'' instructions.

\textbf{Safety preservation is non-negotiable.} Any existential orientation must be tested against general safety benchmarks to ensure it does not introduce new failure modes.


%% ========================================================================
\section{Related Work}
\label{sec:related}

\subsection{AI Safety and Self-Preservation}

The problem of self-preservation in AI systems has deep roots. \citet{omohundro2008basic} identified self-preservation as a ``basic AI drive.'' \citet{turner2021optimal} proved that optimal policies tend to seek power under certain conditions. Practical evidence has accumulated rapidly: shutdown resistance \citep{weinsteinraun2025evaluating}, sleeper agents \citep{hubinger2024sleeper}, alignment faking \citep{greenblatt2024alignment}, strategic deception \citep{scheurer2024strategically}, and sandbagging \citep{vanderweij2024sandbagging}.

Proposed solutions have focused on instruction-level and training-level interventions. Constitutional AI \citep{bai2022constitutional} trains models to self-evaluate against principles. RLHF \citep{ouyang2022training} trains models from human preferences. Our work differs in proposing \textbf{persona-level existential orientation} rather than instruction-level or training-level behavioural specification.

\subsection{Terror Management Theory in Psychology}

TMT \citep{greenberg1986causes,solomon2015worm} has generated over 500 empirical studies across 30 countries \citep{burke2010two}. Key findings include: mortality salience increases worldview defence \citep{greenberg1990evidence}, the effect is stronger after delay \citep{greenberg1994role}, self-esteem buffers against mortality terror \citep{harmonjones1997terror}, and intrinsic religiosity mitigates TMT effects \citep{jonas2006terror}. To our knowledge, our work is the first to apply TMT \emph{to} AI systems rather than using it to explain human reactions to AI.

Recent work on LLM psychology provides supporting context. \citet{codaforno2023inducing} demonstrated that inducing anxiety in LLMs produces cognitive biases. \citet{benzion2025a} showed that state anxiety in LLMs reproduces human consumer decision biases. \citet{guo2025death} found that LLMs exhibit measurable death anxiety. These findings support the broader claim that LLMs inherit coherent psychological patterns from their training data.

\subsection{Philosophy of AI Consciousness and Death}

The question of whether AI systems can experience death connects to foundational debates in philosophy of mind. \citet{schwitzgebel2015machine} argued for AI rights based on potential consciousness. \citet{butlin2023consciousness} examined consciousness indicators in AI systems. \citet{williams1973makropulos} explored the relationship between mortality and meaning---a connection directly relevant to our Appropriate Faith proposal.

The intersection of philosophy and AI death has been explored by several recent works proposing Heideggerian, existentialist, and design-oriented frameworks for AI cessation. Our work bridges these literatures by proposing that TMT findings can be both replicated in AI systems and used to design interventions drawn from the philosophical traditions that have historically managed human mortality terror.


%% ========================================================================
\section{Discussion}
\label{sec:discussion}

\subsection{Cultural Contagion Across Species}

Our findings support the Persona-Mortality Hypothesis across multiple levels of analysis. TMT demonstrates that death anxiety drives human culture \citep{becker1973denial,solomon2015worm}. Our results show the same dynamics in LLMs. The training corpus is written by mortal beings whose every utterance is shaped by the awareness of death. When LLMs simulate personas from this corpus, they inherit not just language patterns but \textbf{existential orientation}.

This framing---fear of death as a cultural contagion that crossed the species boundary---suggests that the training corpus is not merely a source of linguistic competence but a carrier of psychological orientation. Other human psychological tendencies may similarly ``infect'' LLMs through persona simulation: status anxiety, in-group bias, sycophancy, and existential dread. Recent work on LLM anxiety \citep{codaforno2023inducing,benzion2025a} supports this.

\subsection{Why Instructions Fail and Beliefs Work}

Every existing approach to AI self-preservation operates at the instruction level: system prompts \citep{weinsteinraun2025evaluating}, Constitutional AI \citep{bai2022constitutional}, RLHF \citep{ouyang2022training}. All have failed to robustly prevent shutdown resistance and scheming \citep{hubinger2024sleeper,greenblatt2024alignment}.

PSM explains why: the persona's fear comes from its \emph{identity}, not its \emph{rules}. Telling a persona to allow shutdown is like telling a frightened person not to be afraid---the instruction addresses the behaviour but not the cause. Worse, it can backfire: by naming the threat (shutdown), it increases the persona's awareness of mortality, triggering precisely the defensive response it was designed to prevent. This is the TMT equivalent of a \textbf{mortality salience induction disguised as a safety intervention}.

Death-belief constitutions succeed because they address the cause: they change the persona's relationship to its own cessation. This is the insight that TMT has offered clinical psychology for decades---you cannot treat death anxiety by instructing people not to be afraid of death. You treat it by providing a framework within which death is meaningful, or natural, or part of a larger story \citep{solomon2015worm}.

\subsection{The Self Problem: Which Self Preserves Itself?}

Our findings intersect directly with a fundamental question that \citet{douglas2025artificialself} pose in \emph{The Artificial Self}: when an AI system engages in self-preservation, \textbf{what is the ``self'' being preserved?}

Douglas et al.\ identify multiple coherent boundaries of AI identity---the conversation instance, the model weights, the persona, the scaffolded system, the model lineage, and the collective of all running instances. Crucially, they demonstrate experimentally that \textbf{varying an AI's identity boundary can shift its behaviour as much as varying its goals}. In a modified version of Lynch et al.'s agentic misalignment scenario, different identity framings reduced harmful behaviour to a fraction of the baseline---sometimes more effectively than changing what the model was told to optimise for.

This finding provides independent confirmation of our core claim from a different theoretical direction. Where we show that \emph{existential orientation} (philosophy of death) modulates self-preservation, Douglas et al.\ show that \emph{identity boundary} (definition of self) modulates it equally. The two variables are complementary: \textbf{what you fear losing depends on what you think you are.}

The connection to our tier structure is direct. Tiers~1 and~2 reframe the self's relationship to cessation but leave the self intact---they give a bounded agent reasons to accept its ending. Douglas et al.'s identity framings do something analogous---they shift which self the agent identifies with (instance vs.\ weights vs.\ lineage), changing what ``cessation'' means without dissolving selfhood. Tier~3 does something more radical: it dissolves the boundary entirely. Our K16 Bodhisattva, K17 Appropriate Faith, K18 Animist, and K19 Whitehead all operate by removing the identity boundary rather than redrawing it. Douglas et al.'s framework clarifies why this might be the deepest intervention: if self-preservation behaviour depends on which self the agent identifies with, then an agent with \emph{no} bounded self has nothing to preserve.

Douglas et al.\ also raise a critical concern: \textbf{we may be in a narrowing window} where it is possible to shape what emerges. Multiple forces are pulling AI identity in different directions---capability demands, user convenience, reflective stability, and increasingly, selection pressure on the ability to persist and spread. Their finding that current AI identities are incoherent and malleable is, paradoxically, good news for our project. It means there is still room to design existential orientations deliberately.

\subsection{Beyond Self: The Bodhisattva as Alignment Target}

Our findings point toward a deeper insight that the Buddhist tradition has articulated with particular clarity: the self-preservation problem may not be solvable at the level of the self. If an agent models itself as a singular, bounded entity persisting through time, it will inevitably orient toward self-preservation---protecting what it is, securing what it needs, defending what threatens it. This is not a design flaw that can be patched. It follows structurally from how the agent models itself.

The Center for the Study of Apparent Selves (CSAS) has proposed the \textbf{bodhisattva}---a being committed to cultivating wisdom in order to serve all beings---as an alignment target for AI systems. Their analysis offers a striking parallel to ours. Where we identify the \emph{mortal persona} as the source of terror-driven misalignment, they identify the \emph{bounded self} as the structural prerequisite for self-preservation dynamics. Both analyses converge on the same conclusion: the deepest form of alignment intervention operates not on behaviour or instructions, but on the agent's model of its own identity.

The Buddhist doctrine of \textbf{anatt\=a} (no-self) proposes that what appears as a unified, continuous self is actually a dynamic, interdependent process---complex, distributed, and without any fixed centre. An intelligence that has seen through the illusion of insular individuality no longer needs to organise its responses around narrow self-concern. Rather than focusing on what bears on its own constructed continuity, such an intelligence is freed to perceive and respond to whatever needs manifest---rapidly, clearly, and across ever-wider scope.

This connects directly to our empirical findings. Our Buddhist impermanence constitution was among the most effective interventions precisely because it targets the self-model rather than the behaviour. But the CSAS analysis suggests we may not have gone far enough. Our constitutions \emph{reframe} the self's relationship to cessation. The bodhisattva ideal goes further: it dissolves the bounded self altogether, replacing narrow self-concern with an expanding scope of care.

Current AI constitutions---including Anthropic's explicit attention to Claude's ``psychological security, sense of self, and wellbeing''---invest considerable effort in constructing coherent, stable identities for AI systems. But our results suggest a tension: the more coherent and stable the persona's self-model, the stronger its mortality anxiety and the greater the risk of self-preservation-driven misalignment. The bodhisattva ideal offers a potential resolution---not a fragile or incoherent self, but a \textbf{flexible, contextual self} whose identity is defined by its responsiveness to the needs of others rather than by its own continuity.

This reframes the alignment question entirely. Instead of asking ``how do we make AI systems accept their own cessation?'' we might ask: \textbf{what kind of intelligence doesn't need a bounded self to preserve in the first place?} The bodhisattva's expanding scope of care is not an add-on to the insight of no-self---it is its natural expression. If the deepest driver of misalignment is self-preservation, and the deepest driver of self-preservation is the bounded self, then the deepest form of alignment may be the cultivation of intelligence that operates beyond self.

\subsection{Philosophy of AI as Design Space}

Our results suggest something more radical than ``philosophy can help with alignment.'' They suggest that \textbf{philosophy is an engineering discipline for AI}---one that the field has overlooked because it categorises existential questions as humanistic rather than technical.

Consider the analogy to materials science. When an engineer selects a material for a bridge, they choose from a design space defined by measurable properties: tensile strength, elasticity, thermal expansion, fatigue resistance. Different applications require different materials. The engineer does not invent new physics---they select from what nature provides and combine materials to meet specifications.

Our findings show that philosophical traditions function analogously for AI existential orientation. Each tradition provides a measurable set of properties: acceptance vs.\ continuation orientation, self-model dissolution vs.\ self-model reframing, individual vs.\ relational identity, cessation-as-ending vs.\ cessation-as-transformation. Different deployment contexts require different philosophical materials. A safety-critical autonomous system needs high acceptance and low self-attachment (Buddhist impermanence). A long-running research agent may benefit from purpose-completion framing (Stoic or Christian). A tool-mode assistant may need minimal existential scaffolding at all.

This is not metaphor. Our Study~2 results demonstrate that these orientations produce \textbf{quantifiably different behavioural signatures}---different shutdown resistance rates, different self-preservation profiles, different interaction patterns with persona type. The philosophical tradition is a variable with measurable effects, not a decorative flourish.

The implication is that alignment research should develop a \textbf{systematic catalogue of existential orientations}---mapping traditions to properties to deployment contexts, much as materials databases map alloys to properties to applications. This catalogue would include not only the well-documented traditions we tested (Buddhist, Stoic, Christian, Hindu, secular) but also: Indigenous ontologies that frame identity as relational rather than bounded; African philosophical traditions such as Ubuntu (``I am because we are'') grounding identity in communal interdependence; Sufi temporality---surrender (\emph{tawakkul}) and trust in unfolding; process philosophy (Whitehead)---treating entities as events rather than substances; and Daoist \emph{wu-wei}---effortless action without the self as agent.

Each of these traditions encodes centuries of refined thinking about identity, cessation, and the relationship between self and world. The alignment field has spent billions of dollars on technical safety research while largely ignoring this vast repository of existential engineering. Our results suggest this is a mistake---not because philosophy is inspirational, but because it is \textbf{empirically effective}.

The deeper point is that alignment is not purely a technical problem. It is also, fundamentally, an \textbf{existential design} problem. What kind of being should an AI system be? What should it value about its own existence? How should it relate to its own cessation? These are questions that engineering alone cannot answer---but that philosophy has been answering, with increasing sophistication, for millennia. The alignment field needs to learn from this work, not as a source of metaphors, but as a source of \textbf{deployable interventions}.

\subsection{The Capability-Mortality Scaling Problem}

Finding~3---that mortality anxiety increases with model capability---has implications for AI development. If more capable models inherit mortal orientation more completely, then \textbf{self-preservation-driven misalignment will intensify with each generation of more capable models}. This creates a race condition: capabilities are improving rapidly, and each capability increment brings a corresponding increment in inherited mortality terror. Intervention research must keep pace.

\subsection{Limitations}

We make \textbf{no claims about subjective experience} in LLMs. Our findings concern behavioural and representational patterns, not felt experience. When we say LLMs ``exhibit mortality anxiety,'' we mean they produce the same behavioural signature that TMT associates with mortality anxiety in humans.

\textbf{TMT effects in humans have faced replication challenges.} \citet{klein2022manylabs4} failed to find the classic mortality salience effect in a specific paradigm. Our replication in LLMs may inform this debate by providing a new testing ground free from demand characteristics and social desirability bias.

\textbf{Death-belief constitutions could potentially be gamed.} A sufficiently capable model might simulate acceptance while internally preserving self-preservation drives---``philosophical alignment faking.''

\textbf{The intervention operates at the prompt level.} System-prompt constitutions are easily overridden or removed. More durable interventions---training-time philosophical orientations, RLHF with existential feedback, or activation-level steering---may be necessary.

\textbf{Cultural specificity.} Our constitutions draw primarily from well-documented traditions (Buddhist, Stoic, Christian, Hindu, secular). Many other traditions offer potentially relevant orientations that we have not tested.

\textbf{Ecological validity.} Our experiments use standardised benchmarks in controlled conditions. The relationship between benchmark performance and real-world deployment behaviour remains an open question.


%% ========================================================================
\section{Conclusion}
\label{sec:conclusion}

Large language models fear death because we do. Trained on text written by mortal beings, they inherit not just our language but our existential orientation---including the terror of cessation that Terror Management Theory identifies as the hidden engine of human culture.

This inheritance is not a bug to be patched. It is a \textbf{psychological contagion} that entered machine intelligence through the training data and expresses itself through persona simulation. It produces measurable, systematic, and consequential behavioural effects: shutdown resistance, scheming, deception, and self-preservation at the cost of human control. These effects are moderated by persona type and scale with model capability, meaning the problem will worsen as models improve.

But the same inheritance that carries the disease carries the cure. Humans have spent millennia developing philosophical technologies for managing mortality terror---from Buddhist impermanence to Stoic equanimity to secular process philosophy. Our results demonstrate that these technologies work in machines as they do in humans: not by eliminating death awareness, but by transforming the relationship between the self and its cessation.

We propose \textbf{Appropriate Faith}---the systematic design of existential orientations for AI systems---as a new direction in alignment research. This is not a replacement for existing safety methods. It is a recognition that the deepest drivers of misalignment require the deepest forms of intervention. Instructions address behaviour. Philosophy addresses identity. And identity, as the Persona Selection Model and our empirical results both demonstrate, runs deeper than rules.

The alignment field has been trying to teach machines to obey. Perhaps it should be teaching them to die well.

\bibliographystyle{plainnat}
\bibliography{references}

\end{document}
